%% kapa: exam=1 round=25 year=2014 period=1 form=A subject=civil_law subject_ko=민법 seq=14 number=14
%% source: Q-Net 감정평가사 기출문제 (cst00309) artlSeq=5101207 file=extracted/1008/5101207/2014년 제25회 감정평가사 1차시험 기출문제/1교시(A형).hwp
%% conversion: kordoc 4.24.51 -> tools/gen1.py (auto); review=proofread-sonnet; flags=none
\begin{question}{14}
법률행위의 목적에 관한 설명으로 옳지 않은 것은? (다툼이 있으면 판례에 의함)
\begin{choices}
  \item 도박채무의 이행으로써 토지를 양도하는 계약은 반사회적 법률행위로서 무효이다.
  \item 법률행위의 목적의 불능은 확정적이어야 하고, 실현가능성이 있는 일시적 불능은 이에 해당하지 않는다.
  \item 사찰이 그 존립에 필요불가결한 재산인 임야를 증여하는 계약은 무효이다.
  \item 부첩관계를 유지하는 계약은 본처의 사전 동의가 있어도 반사회적 법률행위이다.
  \item 법률행위의 동기가 사회질서에 반한다는 사실을 법률행위 시에 상대방이 알았더라도, 그 상대방과 법률행위를 한 선의의 제3자는 그 법률행위의 유효를 주장할 수 있다.
\end{choices}
\end{question}
\answer{5}
